\subsection{Ablation studies}
\seclabel{ablation}
% from scratch/parameter_count.m:
%   pct of parameters in layers 1-5: 0.06 (3,745,824)
%   pct of parameters in layer 6:    0.65 (37,748,736)
%   pct of parameters in layer 7:    0.29 (16,777,216)
\paragraph{Performance layer-by-layer, without fine-tuning.}
To understand which layers are critical for detection performance, we analyzed results on the VOC 2007 dataset for each of the CNN's last three layers.
Layer \pool{5} was briefly described in \secref{vis}.
The final two layers are summarized below.

Layer \fc{6} is fully connected to \pool{5}. To compute features,
it multiplies a $4096 \times 9216$ weight matrix by the \pool{5} feature map 
(reshaped as a 9216-dimensional vector) and then adds a vector of
biases. 
This intermediate vector is component-wise half-wave rectified ($x \leftarrow \max(0, x)$).

Layer \fc{7} is the final layer of the network. It is implemented
by multiplying the features computed by \fc{6} by a $4096 \times 4096$
weight matrix, and similarly adding a vector of biases and 
applying half-wave rectification.

We start by looking at results from the CNN \emph{without fine-tuning} on PASCAL, \ie all CNN parameters were pre-trained on ILSVRC 2012 only.
Analyzing performance layer-by-layer (\tableref{voc2007} rows 1-3) reveals that
features from \fc{7} generalize worse than 
features from \fc{6}. This means that 29\%, or about 16.8 million,
of the CNN's parameters can be removed without degrading mAP. 
More surprising is that removing \emph{both} \fc{7} and
\fc{6} produces quite good results even though \pool{5}
features are computed using \emph{only 6\%}
of the CNN's parameters. Much of the
CNN's representational power comes from its convolutional layers,
rather than from the much larger densely connected layers. This finding
suggests potential utility in computing a dense feature map,
in the sense of HOG, of an arbitrary-sized image by using only the convolutional layers
of the CNN. This representation would enable experimentation with sliding-window detectors,
including DPM, on top of \pool{5} features.

%\paragraph{Color.}
%To understand how much our system benefits from color (compared to HOG-based
%approaches, which largely ignore it) we tested our pre-trained CNN in a grayscale world.
%Training SVMs on \fc{6} features from grayscale versions of PASCAL images, and 
%testing on grayscale images decreased mAP on VOC 2007 test from 43.4\% to 40.1\%. 

\paragraph{Performance layer-by-layer, with fine-tuning.}

We now look at results from our CNN after having fine-tuned its parameters on VOC 2007 trainval.
The improvement is striking (\tableref{voc2007} rows 4-6):
fine-tuning increases mAP by 8.0 percentage points to 54.2\%. 
The boost from fine-tuning is much larger for \fc{6} and \fc{7} than for \pool{5},
which suggests that the \pool{5} features learned from ImageNet are general and that most of the improvement is gained from learning domain-specific non-linear classifiers on top of them.

\paragraph{Comparison to recent feature learning methods.}
Relatively few feature learning methods have been tried on PASCAL VOC detection.
We look at two recent approaches that build on deformable part models. For reference, we also include results for the standard HOG-based DPM \cite{release5}.

The first DPM feature learning method, DPM ST \cite{lim2013sketch},
augments HOG features with histograms of ``sketch
token'' probabilities. Intuitively, a sketch token is
a tight distribution of contours passing through the center of
an image patch.
Sketch token probabilities are computed at
each pixel by a random forest that was trained to classify
$35 \times 35$ pixel patches into one of 150 sketch tokens or background.

The second method, DPM HSC \cite{HSC}, replaces HOG with histograms of sparse codes (HSC). To compute an HSC, sparse code
activations are solved for at each pixel using a learned dictionary
of 100 $7 \times 7$ pixel (grayscale) atoms. The resulting activations are 
rectified in three ways (full and both half-waves), spatially pooled,
unit $\ell_2$ normalized,
and then power transformed ($x \leftarrow \textrm{sign}(x)|x|^\alpha$).

All R-CNN variants strongly outperform the three DPM baselines (\tableref{voc2007} rows 8-10), including
the two that use feature learning. Compared to the latest version of
DPM, which uses only HOG features, our mAP is more than 20 percentage points higher: 54.2\% \vs 33.7\%---\emph{a 61\% relative improvement}.
The combination of HOG and sketch tokens 
yields 2.5 mAP points over HOG alone, while HSC improves over HOG
by 4 mAP points (when compared internally
to their private DPM baselines---both use non-public implementations of DPM
that underperform the open source version \cite{release5}). These methods achieve mAPs of 29.1\% and 34.3\%, respectively.

\subsection{Network architectures}
\seclabel{netarch}
Most results in this paper use the network architecture from Krizhevsky \etal \cite{krizhevsky2012imagenet}.
However, we have found that the choice of architecture has a large effect on R-CNN detection performance.
In \tableref{voc2007oxfordnet} we show results on VOC 2007 test using the 16-layer deep network recently proposed by Simonyan and Zisserman \cite{vggverydeep}.
This network was one of the top performers in the recent ILSVRC 2014 classification challenge.
The network has a homogeneous structure consisting of 13 layers of $3 \times 3$ convolution kernels,
with five max pooling layers interspersed, and topped with three fully-connected layers.
We refer to this network as ``O-Net'' for OxfordNet and the baseline as ``T-Net'' for TorontoNet.

To use O-Net in R-CNN, we downloaded the publicly available pre-trained network weights for the \texttt{VGG\_ILSVRC\_16\_layers} model from the Caffe Model Zoo.\footnote{\url{https://github.com/BVLC/caffe/wiki/Model-Zoo}}
We then fine-tuned the network using the same protocol as we used for T-Net.
The only difference was to use smaller minibatches (24 examples) as required in order to fit within GPU memory.
The results in \tableref{voc2007oxfordnet} show that R-CNN with O-Net substantially outperforms R-CNN with T-Net, increasing mAP from 58.5\% to 66.0\%.
However there is a considerable drawback in terms of compute time, with the forward pass of O-Net taking roughly 7 times longer than T-Net.
